\chapter{Löser, Faktorisierung und Konvergenz}\label{ch:loeser}

\section{RCWA oder Runge--Kutta?}
\begin{center}
\small
\begin{tabularx}{\linewidth}{>{\raggedright\arraybackslash}p{3.7cm}>{\raggedright\arraybackslash}p{2.6cm}>{\raggedright\arraybackslash}p{3.0cm}Y}
\toprule
Struktur & Empfehlung & Alternative & Begründung \\
\midrule
homogene Schichten, Lamellen, Polygone, Konturen & RCWA & RK & exakt $y$-invariant, ein Eigenproblem je Schicht \\
1D-Profile (Sinus, Dreieck) & RK & RCWA, 40--150 Scheiben & keine Treppe; RK konvergiert mit \key{rk.tol} \\
2D-Profile & RCWA, 20--40 Scheiben & RK für Referenz & RK in 2D sehr teuer \\
CAD-Körper mit schrägen Wänden & RCWA (\key{slicing.rcwa\_slices}) & RK & Bänder je nach Typ \\
Referenzrechnungen & beide & & Übereinstimmung RK $=$ RCWA ist ein starker Test \\
\bottomrule
\end{tabularx}
\end{center}

\subsection{RCWA-Optionen}
\begin{itemize}
  \item \key{rcwa.slices} (Standard 40): Treppenstufen je Profilschicht; je Schicht mit \key{slices} überschreibbar.
        $y$-invariante Schichten brauchen keine Scheiben.
  \item Der Fehler der Treppe fällt etwa wie $1/\text{slices}$. Beim Al-Sinus V4 (TM) ergeben 60, 150 und 300
        Scheiben 0,853, 0,865 und 0,869 -- der Grenzwert ist 0,8696. Für Metallprofile in TM sind also viele Scheiben
        nötig; dort ist RK meist schneller.
\end{itemize}

\subsection{RK-Optionen}
\begin{itemize}
  \item \key{rk.tableau}: \val{gauss2} (Standard, Ordnung 4), \val{gauss3} (Ordnung 6, weniger Schritte bei hoher
        Genauigkeit), \val{radau3} (L-stabil, für steife Fälle).
  \item \key{rk.tol} (Standard $10^{-8}$): lokale Toleranz der Schrittweitensteuerung durch Schrittverdopplung. Für
        Übersichtsrechnungen genügt $10^{-5}$; die Laufzeit sinkt dann deutlich.
  \item \key{growth\_limit}: wann eine neue Stabilitätsscheibe beginnt (0 = automatisch, in fp64 $10^8$). Nur bei
        numerischen Problemen ändern.
\end{itemize}
Die Diagnose (\texttt{rk\_steps}) zeigt, wie viele Schritte gebraucht wurden.

\section{Faktorisierung}\label{sec:faktorisierung}
Die Faktorisierung legt fest, wie das Produkt aus Permittivität und Feld in Fourierreihen gebildet wird. Sie ist der
entscheidende Faktor für die Konvergenzgeschwindigkeit bei TM und bei Metallen.
\begin{table}[h]
\centering
\small
\begin{tabularx}{\linewidth}{>{\raggedright\arraybackslash}p{3.3cm}>{\raggedright\arraybackslash}p{3.0cm}X}
\toprule
Schichttyp & \val{auto} wählt & zulässig \\
\midrule
\val{lamellar1d} isotrop & \val{li\_lamellar} & \val{laurent}, \val{li\_lamellar} (die FFF-Namen wirken hier wie \val{li\_lamellar}) \\
\val{lamellar1d} anisotrop & \val{FFF\_A} & \val{FFF\_A}, \val{FFF\_B}, \val{laurent} \\
\val{profile1d} isotrop & \val{FFF\_rig} (Sinus), \val{FFF\_sym} (Dreieck, Polygon) & \val{laurent}, \val{li\_lamellar},
\val{FFF\_iso}, \val{FFF\_sym}, \val{FFF\_rig} \\
\val{profile1d} anisotrop & \val{FFF\_A} & \val{FFF\_A}, \val{FFF\_B}, \val{laurent} \\
\val{polygon2d} & \val{FFF\_sym} / \val{FFF\_B} (aniso.) & \val{FFF\_sym}, \val{FFF\_iso}, \val{FFF\_B}, \val{laurent} \\
\val{contour2d} & \val{FFF\_sym} / \val{FFF\_B} (aniso.) & \val{FFF\_sym}, \val{FFF\_iso}, \val{FFF\_B}, \val{laurent} \\
\val{profile2d} & \val{FFF\_rig} / \val{FFF\_A} (aniso.) & isotrop: \val{laurent}, \val{FFF\_iso}, \val{FFF\_sym}, \val{FFF\_rig}, \val{FFF\_A}, \val{FFF\_B}; anisotrop: \val{laurent}, \val{FFF\_A}, \val{FFF\_B} \\
\bottomrule
\end{tabularx}
\caption{Faktorisierungen je Schichttyp. Groß- und Kleinschreibung der FFF-Namen ist egal.}
\label{tab:faktorisierung}
\end{table}

\begin{description}
  \item[\val{laurent}] Naive Produktregel. Korrekt im Grenzwert, aber bei TM und Metallen sehr langsam. Nur für
        Vergleiche. Bei RCWA-Lamellen ist sie dafür exakt energieerhaltend.
  \item[\val{li\_lamellar}] Li-Regeln für Sprünge in $x$: Inverse Regel für die Normalkomponente. Standard für
        Lamellen.
  \item[\val{FFF\_rig}] FFF mit glattem Normalenfeld; beste Wahl für glatte Profile (Sinus).
  \item[\val{FFF\_sym}] symmetrisierte FFF; robust bei Profilen mit Ecken und bei 2D-Querschnitten.
  \item[\val{FFF\_iso}] Formel aus dem Buch von Nevière/Popov. \textbf{Versagt bei schwach absorbierenden Metallen}
        (Energiebilanz $>1$); \FG{} gibt dann eine Warnung aus.
  \item[\val{FFF\_A}, \val{FFF\_B}] Varianten für anisotrope Materialien (Profile bzw.\ Querschnitte).
\end{description}
\textbf{Faustregel:} Lassen Sie \val{auto} stehen. Ändern Sie die Faktorisierung nur für Vergleiche oder wenn die
Fehlersuche (Kapitel~\ref{ch:probleme}) es nahelegt.

\section{Konvergenz prüfen}\label{sec:konvergenz}
\subsection{Vorgehen}
\begin{enumerate}
  \item Rechnen Sie mit einem Startwert von $N$ (Richtwerte unten).
  \item Verdoppeln Sie $N$ (in 2D: erhöhen Sie $N_x$ und $N_z$ um 50\,\%) und rechnen Sie erneut.
  \item Ändern sich die Effizienzen, die Sie interessieren, um weniger als Ihre Genauigkeitsanforderung, ist das
        Ergebnis brauchbar. Sonst weiter erhöhen.
  \item Bei RCWA-Profilen ebenso die Scheibenzahl, bei RK die Toleranz prüfen.
\end{enumerate}
Am bequemsten geht das mit einem Sweep über \val{/truncation/Nx}:
\begin{lstlisting}[language=json]
"sweep": { "parameter": "/truncation/Nx", "values": [10, 20, 40, 80] }
\end{lstlisting}

\subsection{Richtwerte für $N$}
\begin{center}
\small
\begin{tabular}{lll}
\toprule
Struktur & Start & typisch ausreichend \\
\midrule
1D-Dielektrikum & 10 & 15--20 \\
1D-Metall, TE & 15 & 20--30 \\
1D-Metall, TM & 20 & 30--40 (mit ASR 20) \\
2D-Dielektrikum & 4 & 5--8 je Richtung \\
2D-Metall (Lochgitter) & 5 & 8--12 je Richtung \\
\bottomrule
\end{tabular}
\end{center}

\subsection{Ein Beispiel: Goldlamellen in TM}
Goldlamellen ($\nu=0{,}2+3{,}4i$, Periode 1\um, Breite 0,5\um, Höhe 0,3\um) auf Gold, $\lambda=0{,}6328$\um,
$\varphi=20^\circ$, TM. Das ist der Referenzfall V19; der Grenzwert aus einer Extrapolation bis $N=800$ ist
$e_{-1}=0{,}013591$, $e_0=0{,}79435$. Gerechnet mit \FG{} 0.18:
\begin{center}
\small
\begin{tabular}{rcccccc}
\toprule
& \multicolumn{2}{c}{\val{laurent}} & \multicolumn{2}{c}{\val{li\_lamellar} (= \val{auto})} & \multicolumn{2}{c}{\val{li\_lamellar} + ASR $\eta=0{,}95$} \\
\cmidrule(lr){2-3}\cmidrule(lr){4-5}\cmidrule(lr){6-7}
$N$ & $e_{-1}$ & $e_0$ & $e_{-1}$ & $e_0$ & $e_{-1}$ & $e_0$ \\
\midrule
10 & 0,01770 & 0,7314 & 0,01282 & 0,7938 & 0,01374 & 0,79422 \\
20 & 0,02392 & 0,6804 & 0,01313 & 0,7922 & 0,01358 & 0,79421 \\
40 & 0,01405 & 0,7007 & 0,01334 & 0,7930 & 0,01358 & 0,79427 \\
80 & 0,00983 & 0,6829 & 0,01346 & 0,7936 & 0,01358 & 0,79429 \\
\midrule
Grenzwert & 0,013591 & 0,79435 & 0,013591 & 0,79435 & 0,013591 & 0,79435 \\
\bottomrule
\end{tabular}
\end{center}
Was man daran sieht:
\begin{itemize}
  \item Mit \val{laurent} springt das Ergebnis auch bei $N=80$ noch um 30\,\% -- unbrauchbar.
  \item Mit den Li-Regeln (Standard) stimmt $e_0$ schon bei $N=10$ auf 0,1\,\%; $e_{-1}$ konvergiert langsamer.
  \item Mit der adaptiven räumlichen Auflösung (ASR, Abschnitt~\ref{sec:asr}) ist $e_{-1}$ ab $N=20$ auf vier Stellen
        richtig.
\end{itemize}

\subsection{Resonanzen}
In der Nähe von Resonanzen (Oberflächenplasmonen, geführte Moden, Rayleigh-Anomalien) verschiebt sich die
\emph{Lage} eines Peaks mit $N$. Ein Vergleich bei einer festen Wellenlänge kann dann große Unterschiede zeigen,
obwohl die Rechnung im Wesentlichen stimmt. Vergleichen Sie in solchen Fällen ganze Spektren: Lage und Höhe des
Peaks müssen konvergieren (Referenzfall V10, Silber-Lochgitter).

\section{Gütekriterien einer Rechnung}\label{sec:guete}
\begin{itemize}
  \item \textbf{Energiebilanz}: Bei verlustfreien Strukturen muss \texttt{sum R+T} $=1$ sein -- RCWA mit Laurent
        exakt, FFF-Varianten auf $10^{-5}$ bis $5\cdot10^{-4}$. Größere Abweichungen zeigen ein zu kleines $N$ oder eine
        ungeeignete Faktorisierung.
  \item \textbf{Absorption $\ge0$}: Negative Absorption bei verlustbehafteten Strukturen ist ein Warnsignal.
  \item \textbf{RK $=$ RCWA}: Für Lamellen und Polygone müssen beide Löser übereinstimmen (Toleranz der RK-Rechnung).
  \item \textbf{Reziprozität}: Vertauschen von Einfall und Beugungsrichtung muss dieselbe Effizienz liefern (mit $N$
        konvergierend).
  \item \textbf{Grenzfälle}: Struktur mit Tiefe 0 $=$ Fresnel; $z$-invariante 2D-Struktur $=$ 1D-Rechnung.
\end{itemize}
